%
% sesion4.tex
% 
% Copyright 2017 Rony J. Letona <zronyj@gmail.com>
% 
% This program is free software; you can redistribute it and/or modify
% it under the terms of the GNU General Public License as published by
% the Free Software Foundation; either version 2 of the License, or
% any later version.
% 
% This program is distributed in the hope that it will be useful,
% but WITHOUT ANY WARRANTY; without even the implied warranty of
% MERCHANTABILITY or FITNESS FOR A PARTICULAR PURPOSE.  See the
% GNU General Public License for more details.
% 
% You should have received a copy of the GNU General Public License
% along with this program; if not, write to the Free Software
% Foundation, Inc., 51 Franklin Street, Fifth Floor, Boston,
% MA 02110-1301, USA.
%

\documentclass[10pt,letterpaper]{article}
\usepackage[latin1]{inputenc}
\usepackage[spanish]{babel}
\usepackage{graphicx}
\usepackage{hyperref}
\usepackage{amsmath}
\usepackage{amsfonts}
\usepackage{amssymb}
\usepackage{color}
\usepackage{float}
\usepackage{upquote}
\usepackage[left=2cm,right=2cm,top=2cm,bottom=2cm]{geometry}
\author{Rony J. Letona}
\title{Taller de Computaci\'on Cient\'ifica para Ciencias Qu\'imicas: Sesi\'on 4}
\definecolor{light-gray}{gray}{0.90}

\newcommand{\tab}[1]{\hspace{.2\textwidth}\rlap{#1}}

\newcommand{\inlinecode}[1]{
\colorbox{light-gray}{\texttt{#1}}
}

\newsavebox{\selvestebox}
\newenvironment{Code}
{
\begin{lrbox}{\selvestebox}%
\begin{minipage}{\dimexpr\columnwidth-2\fboxsep\relax}
\fontfamily{\ttdefault}\selectfont
}
{\end{minipage}\end{lrbox}%
\begin{center}
\colorbox{light-gray}{\usebox{\selvestebox}}
\end{center}
}

\newcommand{\Picture}[1]
{
	\begin{figure}[H]
	\begin{flushleft}
	\includegraphics[width=\columnwidth]{#1}
	\end{flushleft}
	\end{figure}
}

\begin{document}
\maketitle

\section{Principios de Programaci\'on para las Ciencias 1}
Para muchos de nosotros, ver la palabra \emph{programaci\'on} ya es motivo suficiente para decir ``Est\'an locos si creen que voy a hacer esto otra vez''. Probablemente ya hemos intentado darle instrucciones a un ordenador en el pasado (posiblemente en el colegio) y no ha sido nada bonito. Sin embargo, ahora no va a ser como en ese entonces.\\

En el taller hemos ido viendo las cosas lento y paso a paso. No nos van a pedir proyectos ni ejercicios traum\'aticos con los que tendremos pesadillas que ni el psic\'ologo querr\'a escuchar. No, esta vez ser\'a mucho m\'as f\'acil y mucho m\'as interesante. Si el programa no corre o se rompe, tranquilos, no pasa nada. Pero para tener todo claro, comencemos con la pregunta de muchos: Para qu\'e programaci\'on?\\

Mucho del software escrito para fines cient\'ificos resulta ser, generalmente, muy espec\'ifico. Este hace una tarea \textbf{muy} bien, pero no permite hacer algunas otras cosas que nos interesan. Sin embargo, mucho del software permite que se le escriba plugins, extensiones o paquetes con los que podemos extender la funcionalidad del programa original. Adem\'as de esto, el poder programar nuestras propias rutinas o pruebas nos da la capacidad de analizar mejor nuestros datos, calcular propiedades espec\'ificas (o a veces nuevas) y llevar a cabo procesos en el orden que nos interesa. En todas las ciencias, esto resulta ser muy pr\'actico.\\

Otra pregunta que tambi\'en puede surgirnos es: Por qu\'e Python como lenguaje de programaci\'on? Hasta hace una d\'ecada m\'as o menos, el lenguaje utilizado para hacer c\'alculos y rutinas en ciencia era Fortran. Este lenguaje, sin embargo, se quedaba atr\'as en lo que se refiere a versatilidad. Solo serv\'ia para c\'alculos y no permit\'ia m\'as. Luego est\'a Pascal o Delphi. Muchos aprendimos a programar en \'el y hasta la fecha se sigue usando para aprender a programar, pero solo para eso; no se utiliza casi nunca en el mundo real. Visual Basic es otro de aquellos lenguajes que algunos aprendimos. Funciona bien casi solo en ambientes Microsoft Windows, lo cual nos limita bastante. Tambi\'en est\'an las alternativas como C y C++. Estos dos, seg\'un muchos educadores, nos confunden m\'as a los estudiantes con la sintaxis de lo que nos pueden servir para aprender la l\'ogica de programaci\'on. El caso es similar con Java, aunque este tiene una ventaja que los anteriores no: un programa corre sobre cualquier sistema operativo; no hace falta compilarlo para cada caso. JavaScript es el famoso lenguaje utilizado en p\'aginas web y, con el nacimiento de Node.JS, ahora se puede utilizar para programas formales que no necesitan un navegador. El problema es que todav\'ia no cuenta con tantos paquetes para extender su funcionalidad. Finalmente, desde el 7 de julio de 2014, es un hecho que la mayor\'ia de universidades de renombre en Canad\'a, Estados Unidos y varias en Europa utilizan Python como el lenguaje a aprender en sus cursos b\'asicos de programaci\'on para todas las carreras.\\

Las ventajas que tiene Python sobre los dem\'as lenguajes son todas las antes mencionadas: versatilidad, cantidad de paquetes, corre en cualquier plataforma, es f\'acil de aprender y es bastante universal. Adem\'as, existe bastante documentaci\'on para \'el y mucho del software para Computai\'on Cient\'ifica est\'a parcialmente escrito en \'el. Lo \'unico en lo que tiene una clara desventaja es en su velocidad. Sin embargo, para efectos pr\'acticos, esto no importa tanto ya que no trabajaremos con \emph{demasiadas} operaciones y nuestros procesadores ya son capaces de suplir algunas de las deficiencias en velocidad de este lenguaje. Si dese\'aramos trabajar con un lenguaje m\'as r\'apido y que funcione en todas partes, Java es quiz\'a la siguiente mejor opci\'on, aunque la sintaxis nos exigir\'a ser m\'as rigurosos a la hora de escribir nuestros programas.\\

Para poder responder mejor a todas nuestras dudas, comencemos con los ejercicios y vamos a irnos dando cuenta c\'omo es que con pocos comandos e informaci\'on, logramos hacer desde peque\~nos c\'alculos hasta an\'alisis de datos y predicci\'on de propiedades.\\

\subsection{Operaciones B\'asicas}
Vamos a comenzar con lo b\'asico. Vamos a abrir nuestra l\'inea de comando, vamos a escribir \inlinecode{python} en ella y vamos a presionar enter. Inmediatamente veremos la interfaz lista para comenzar: \inlinecode{$>>>$} Una cosa que tenemos que tomar en cuenta es que vamos a tener que ir re-aprendiendo a hacer muchas cosas al aprender a programar. Se dice que al aprender a programar, la forma de pensar de una persona cambia un poco. Algo as\'i como cuando aprendemos a ser rigurosos en un an\'alisis.\\

La idea con esta primera parte ser\'a entonces comenzar haciendo ejercicios \textbf{muy} simples. Desde aritm\'etica sencilla hasta almacenar valores en variables. Estas cosas tan sencillas se pueden hacer en calculadoras, pero una cosa que debemos recordar es que un ordenador no es m\'as que eso: una calculadora muy sofisticada. Pero si aprendemos a usarla bien, no habr\'a operaci\'on que no podamos hacer. Comencemos!

\subsubsection{Aritm\'etica}
Para comenzar vamos a tomar la operaci\'on m\'as b\'asica que aprendimos: sumar. C\'omo se hace esto? Pues es sencillo, hacemos algo as\'i: \inlinecode{4\ +\ 5} y presionamos enter. Como nos damos cuenta, esto es relativamente f\'acil. Es una buena pr\'actica en programaci\'on dejar espacios a los lados del signo de la operaci\'on; en este caso la suma \inlinecode{+}\\

Intentemos ahora realizar una resta. Atrevidamente podremos pensar de una vez en cosas como ``Y qu\'e pasa si la resta tiene un resultado menor a 0?'' Pues la manera de averiguarlo es probando: \inlinecode{8\ -\ 13} El resultado era lo que esper\'abamos? Sigamos probando.\\

Ahora multipliquemos. Esto tampoco se ve dif\'icil: \inlinecode{14\ *\ 86} De hecho, podemos notar que la velocidad de respuesta es igual que la de una calculadora. Si as\'i lo deseamos, podemos hacer m\'as pruebas que las que aqu\'i estamos poniendo. Por ahora, todo va bien.\\

Llegamos a la divisi\'on. Intentemos lo siguiente entonces: \inlinecode{16\ /\ 3} Obtuvimos una respuesta? Perfecto! La respuesta era la que esper\'abamos? No? Por qu\'e no? Por qu\'e ser\'a que la respuesta no sali\'o como esper\'abamos? Despu\'es de todo esto solo fue una simple divisi\'on de enteros. Ah! All\'i est\'a la clave! Fue una divisi\'on de \emph{enteros}. Eso significa que el ordenador todav\'ia no sabe que estamos trabajando con decimales, as\'i que no los toma en cuenta. C\'omo hacer que esto funcione entonces? Pues indic\'andole que debe de considerar que los n\'umeros pueden tener decimales, as\'i: \inlinecode{16.0\ /\ 3}\\

Ya vimos las 4 operaciones b\'asicas que conocemos. Ahora, recordando algo que vimos en la sesi\'on 2, vamos a ver c\'omo se calculan los residuos de una divisi\'on de enteros mediante la operaci\'on \emph{m\'odulo}: \inlinecode{\%} Si intentamos una divisi\'on de n\'umeros enteros, como \inlinecode{18\ /\ 7}, vamos a obtener \inlinecode{2} como resultado. Pero sabemos que \inlinecode{7\ *\ 2} no es \inlinecode{18} De la divisi\'on anterior obten\'iamos un residuo que podemos calcular as\'i: \inlinecode{18\ \%\ 7} A primera vista esto no se ve de mucha utilidad, pero al trabajar con listas, arreglos o conjuntos de objetos vamos a ver que nos va a servir.\\

Finalmente, vamos a ver una operaci\'on que no es tan b\'asica, pero nos puede servir mucho. Para evitar estar multiplicando muchas veces el mismo n\'umero, en matem\'atica se crearon las potencias. Una operaci\'on que se ve sencilla, como $3^{5}$ o $2^{7}$ nos toman un momento al hacerlas en la mente, pero resulta que nuestro ordenador lo puede hacer en mil\'esimas de segundo. Intentemos esas dos operaciones! Para la primera, ingresamos \inlinecode{3\ **\ 5} y observamos el resultado. Efectivamente obtuvimos lo que dese\'abamos. Ahora intentemos la segunda: \inlinecode{2\ **\ 7} Tambi\'en obtenemos algo que nos satisface. Pero, qu\'e pasa si intentamos calcular raices o potencias con decimales? Intentemos dos ejemplos: \inlinecode{3.14159265359\ **\ (1/2.0)} y \inlinecode{5\ **\ 2.71828182846} Muy interesante! Con esto notamos que podemos hacer muchas m\'as operaciones. Tomemos nota y enumeremos las operaciones que podemos llevar a cabo con todo lo que acabamos de aprender.

\subsubsection{Variables}
Por ahora vamos viendo que el ordenador se comporta como una plataforma que nos permite hacer operaciones matem\'aticas. Por eso, a la hora de escuchar la palabra \emph{variable}, es natural que pensemos que estas son como las que incluimos en las funciones (e.g. \emph{x} en $ f \left( x \right) = x^{2} - x - 1 $). Sin embargo, en el ambiente de un ordenador esto es un poco diferente. Las variables que nosotros conocemos pueden no tener un valor al inicio, como acabamos de ver. Pero al trabajar en un ambiente digital, este requiere n\'umeros, valores y algo s\'olido con qu\'e trabajar. Por ello, podemos pensar que las variables aqu\'i ser\'an solo un lugar d\'onde almacenar datos, lo cual no es nada lejano a la realidad.\\

Un ordenador posee un dispositivo llamado memoria. En ella se almacenan datos de manera vol\'atil (i.e. se borran al apagar el sistema). Desde que arrancamos nuestro ordenador hasta que lo apagamos, la memoria (RAM) est\'a alojando informaci\'on que permite que trabajemos en el momento. Comprendemos ahora que \emph{ese} es aquel lugar en donde est\'an nuestros documentos antes de que guardemos cambios. Por eso si nos quedamos sin energ\'ia el\'ectrica, todo se pierde. Aparte es el espacio en el disco duro, el cual nos da la capacidad de almacenar datos de manera permanente. Pero en esta secci\'on vamos a ver un concepto r\'apido que nos servir\'a para comprender c\'omo funcionan las variables y otras cosas m\'as adelante.\\

Volviendo a las variables, estas son entonces un espacio en memoria. Algo como una peque\~na caja $ \square $ en donde se puede almacenar un dato particular. Podemos sacarlo de la caja? Claro! Tambi\'en podemos colocar algo diferente all\'i despu\'es. Como espacio en memoria, puede almacenar lo que sea. Eso s\'i, no pueden haber dos datos en la misma caja, y por ende, no pueden haber dos datos almacenados en la misma variable.\\

Pasemos, pues, a almacenar algunos datos en variables. La forma de hacer esto es escribendo algo as\'i: \inlinecode{a = 5} Presionamos enter y ... nada pas\'o? De hecho, si no ha salido ning\'un error, todo est\'a bien. Lo que pasa es que no hay nada que mostrar. El dato fue almacenado y ya. Si queremos ver lo que hay dentro de la variable, podemos escribir lo siguiente: \inlinecode{print(a)} Esto nos mostrar\'a o \emph{imprimir\'a} en pantalla aquello que hay en la variable \inlinecode{a} Debe de llamarse \emph{a} o \emph{b} o \emph{c}? No, puede llamarse como nosotros querramos! Lo \'unico que hay que tomar en cuenta en una variable es que su nombre no puede comenzar con un n\'umero, y esta no puede tener caracteres con tildes, di\'eresis, etc. Tampoco e\~nes. Otra cosa importante es que no podemos ponerle dos nombres a una variable. Eso significa que no puede haber espacios en el nombre de una variable. Con esto en mente, intentemos otra cosa: \inlinecode{mi\_segunda\_variable = 1.61803398875} Como podemos ver, esto tambi\'en se vale. Si despu\'es imprimimos el valor de \inlinecode{mi\_segunda\_variable} veremos que todo funciona bien. Intentemos algo m\'as interesante con esto.\\

Vamos a intentar transformar una temperatura de grados Farenheit a Kelvin. Esto implica transformar primero la temperatura de grados Farenheit a grados Celcius, y luego a Kelvin. Para ello, hay 3 n\'umeros importantes que debemos recordar: una pendiente \inlinecode{m = 5.0/9}, un desplacamiento en el eje \emph{x} \inlinecode{d = 32} y un desplazamiento en el eje \emph{y} \inlinecode{b = 273.15} Ahora veamos, la f\'ormula ir\'ia algo as\'i:

\begin{itemize}
\item De $^o F$ a $^o C$:\hspace*{1cm} $T_{C} = \frac{5}{9} \left( T_{F} - 32 \right) = m \left( T_{F} - d \right)$
\item De $^o C$ a $K$:\hspace*{1cm} $T_{K} = T_{C} + 273.15 = T_{C} + b$
\end{itemize}

Si sustituimos una dentro de la otra, nos queda una f\'ormula as\'i: $T_{K} = m \left( T_{F} - d \right) + b$. Ahora, si nos dicen que la temperatura hoy es de $73^o F$ lo que hacemos nosotros, habiendo almacenado las variables anteriores, es esto: \inlinecode{m * (73 - d) + b} El resultado sale de inmediato y nos complace saber que lo hicimos bien.\\

Perfecto! Ahora ya sabemos usar variables. Como \'ultimo ejercicio antes de la siguiente secci\'on, intentemos convertir la siguiente temperatura de Kelvin a grados Farenheit y luego a Celsius: $298.45 K$.

\subsection{Tipos de Datos}
En qu\'imica, cuando comenzamos a ver de qu\'e est\'a conformada la materia, aprendemos que esta tiene diferentes partes: mol\'eculas, \'atomos, iones, electrones, protones, etc. Aqu\'i sucede algo similar. La informaci\'on en nuestro ordenador se puede descomponer en diferentes partes o \emph{tipos de datos}. Por ahora solo hemos conocido a los n\'umeros enteros y aquellos con punto decimal, pero ahora vamos a comenzar desde un poco m\'as abajo.\\

Ya es cultura general que la informaci\'on en nuestros ordenadores se almacena en forma de unos y ceros, pero esto no es solo cuesti\'on de que si deseamos vamos a ver la informaci\'on codificada as\'i. Lo que tenemos que tomar en cuenta es que estos n\'umeros son c\'odigo binario que codifica a n\'umeros enteros. Estos, a su vez, representan a veces letras y valores num\'ericos o l\'ogicos. Comencemos viendo estos \'ultimos.\\

Al almacenar alg\'un dato en una variable, podemos averiguar qu\'e tipo de dato es el que contiene la variable por medio de un comando sencillo: \inlinecode{type()} Al momento de querer saber de qu\'e tipo de dato se trata, solo aplicamos la funci\'on \textbf{type} a una variable de la manera siguiente: \inlinecode{type(variable)} Conforme vayamos avanzando en este tema, iremos viendo c\'omo es que Python le llama a cada tipo de dato.

\subsubsection{Booleanos}
Los valores l\'ogicos son solo dos: \emph{verdadero} y \emph{falso}. Si recordamos nuestros cursos de l\'ogica y filosof\'ia, recordamos aquellas tablas de verdad que se nos pon\'ia a hacer para revisar el valor de una proposici\'on. A pesar de que no vamos a hacer tablas de verdad, es bueno recordar c\'omo funcionaban, puesto que aqu\'i vamos a usar muchas de las cosas que all\'i se usaban. Pero bueno, comencemos a ver c\'omo funcionan los valores l\'ogicos.\\

Como ya hab\'iamos mencionado, existen solo dos tipos de valor l\'ogico a los que denominamos \emph{booleanos} por el inventor de estos. Para \textbf{verdadero} vamos a tener \inlinecode{True} y para \textbf{falso} \inlinecode{False}. Estos los podemos colocar en nuestro IDE y vamos a ver que toman un color particular poniendo en evidencia de que se trata de un valor particular. Pero solos estos no hacen mucho.\\

Algo que siempre se ve de la mano con los valores, son los operadores l\'ogicos. Cuando los conocemos, nos los presentan con nombres un poco complicados como \emph{conjunci\'on}, \emph{disyunci\'on} y \emph{negaci\'on}. En este caso en particular, vamos a verlos en acci\'on de otra forma. Cuando deseamos que dos cosas se cumplan para obtener un resultado verdadero usamos el operador \inlinecode{and}. Cuando deseamos que al menos una de las dos cosas se cumpla para obtener un resultado verdadero usamos \inlinecode{or}. Y cuando deseamos negar un valor en particular, utilizamos \inlinecode{not}. Esto nos va a servir m\'as adelante, cuando querramos tomar decisiones en un programa o una rutina. Para mientras, vamos a ver de d\'onde salen los booleanos.\\

Para tomar una decisi\'on, generalmente revisamos una comparaci\'on de cierto tipo: es esto igual a aquello, es esto m\'as grande que aquello, es esto diferente a aquello, etc. Este tipo de comparaciones se pueden hacer en nuestro ordenador de la siguiente manera:

\begin{enumerate}
\item Si deseamos revisar si dos cosas son iguales: \inlinecode{a == b}
\item Si deseamos revisar si esas dos cosas son diferentes: \inlinecode{a != b}
\item Si deseamos revisar si una cosa es mayor a la otra: \inlinecode{a >\ b}
\item Si deseamos revisar si esa cosa es menor a la otra: \inlinecode{a <\ b}
\item Si deseamos revisar si una cosa es mayor o igual a la otra: \inlinecode{a >= b}
\item Si deseamos revisar si esa cosa es menor o igual a la otra: \inlinecode{a <= b}
\end{enumerate}

Estas operaciones, al igual que las operaciones matem\'aticas, pueden llevarse a cabo combinadas con m\'as operaciones l\'ogicas. Con esto ya podemos escribir cosas como \inlinecode{(a == 0) and (b != 10)} Claro que para hacer todo eso necesitamos que \textbf{a} y \textbf{b} tengan alg\'un valor, pero ahora sabemos que podemos comparar valores de variables y trabajar con ellos. Y antes de pasar a la siguiente secci\'on, conviene decir que Python reconoce este tipo de dato como \inlinecode{bool}

\subsubsection{N\'umeros}
Este tipo de datos ya lo conocemos ... parcialmente. Sabemos que existen los \textit{enteros} (e.g. 1, -4, 7, 0, 23, ...), los cuales en Python son representados por \inlinecode{int} y los utilizamos generalmente para enumerar cosas. Los enteros llegan hasta donde la memoria de nuestro ordenador nos deje. S\'i, resulta que estos tienen un l\'imite, porque despu\'es de cierto n\'umero, la cifra es tan alta que una cajita de la memoria ya no es suficiente.\\

Despu\'es de esos n\'umeros comienzan los n\'umeros \textbf{largos} (e.g. 9223372036854775808L, -9223372036854775809L, ...). Estos son enteros MUY grandes y si ponemos atenci\'on, tienen una L al final denotando su extra\~no tipo. En Python son representados por \inlinecode{long} y solo nos sirven para c\'alculos con cifras muy grandes. Fuera de eso, cualquier \inlinecode{int} nos sirve para tareas cotidianas.\\

Luego tenemos a un tipo de n\'umeros mucho m\'as conocidos: los de \textit{punto flotante}. Los \emph{qu\'e}? Los n\'umeros de punto flotante. Estos n\'umeros son aquellos que tienen decimales (e.g. 0.333333333333, 3.141592653589, ...). Se les llama de punto flotante porque lo que nuestro  ordenador est\'a haciendo realmente, es guardar el n\'umero sin el punto decimal y guardar la posici\'on del punto en otro espacio en memoria. Al llamar a un n\'umero con decimales, nuestro ordenador va por el n\'umero, y coloca el punto decimal, como flotando, en la posici\'on indicada. Un dato importante a tomar en cuenta es que Python solo trabaja con 12 decimales, y este tipo de dato es representado por \inlinecode{float}\\

Finalmente, tenemos a un tipo de n\'umero con el que uno no se topa muy seguido: los \textit{n\'umeros complejos}\\(e.g. 3 + 4i, -7 - 2i, ...). S\'i, Python puede manejar n\'umeros complejos represent\'andolos como \inlinecode{complex}, solo que en vez de utilizar una \textbf{\^{i}}, se utiliza una \textbf{j} para denotar la parte imaginaria del n\'umero. El uso de los n\'umeros complejos no es algo que se haga muy seguido en el campo de la qu\'imica. De hecho, solo se utilizan casi para c\'alculos cu\'anticos, y a\'un esos los evitamos muchas veces.

\subsubsection{Cadenas}
Despu\'es de ver que podemos almacenar valores de \emph{verdadero}, \emph{falso} y varios n\'umeros, es solo natural que alguien pregunte: ``Y c\'omo se trabaja con texto?'' El texto es un tipo de dato muy particular. Este es una estructura que consta de muchas peque\~nas partes: las letras. Por ser una estructura alargada y que consta de elementos que la unen, a los fragmentos de texto se les denomin\'o cadenas. En ingl\'es se les conoce como \emph{strings}, por pensar en el mismo principio, pero con un cord\'on.\\

Las cadenas se representan en Python como \inlinecode{str}, y para crearlas, colocamos el texto que deseamos entre comillas o ap\'ostrofes. Intentemos con un ejemplo: Vamos a guardar el texto \textit{Hola mundo!} en una variable. Para ello vamos a ingresar \inlinecode{mi\_variable = "Hola mundo!"} Solamente! No hay que hacer nada extra\~no o diferente. Ahora, si deseamos mostrar lo que guardamos, solo imprimimos lo que hay en la variable: \inlinecode{print(mi\_variable)}\\

Perfecto, ya podemos guardar texto. Ahora veamos qu\'e m\'as se puede hacer con \'el, pero intentando entrar ya un poco en qu\'imica. Asumamos, pues, que podemos representar a un hidrocarburo como una estructura de Lewis, ignorando sus hidr\'ogenos y sin representar a los enlaces. Si estamos hablando de una cadena de 6 \'atomos de carbono, el resultado ser\'ia as\'i: \inlinecode{hexano = "\ \hspace{-2mm}CCCCCC"} Como podemos ver, estamos representando a un hidrocarburo con una simple cadena de texto. Qu\'e pasa si deseamos saber cu\'antas letras tiene la cadena? Pues utilizamos una funci\'on para ello: \inlinecode{len(hexano)}\\

Una cosa importante que tambi\'en es conveniente saber es que las cadenas se pueden sumar. Lo que se est\'a haciendo realmente, es concatenarlas. Si de nuestra mol\'ecula anterior queremos hacer un octano, solo nos hacen falta dos carbonos. Entonces, podemos decir que \inlinecode{octano = hexano + "\ \hspace{-2mm}CC"} Y de esta manera ya tenemos una nueva cadena con nuevos elementos.\\

Intentemos crear otra mol\'ecula: \inlinecode{dietileter = "\ \hspace{-2mm}CCOCC"} En esta ya incluimos otro \'atomo diferente del carbono. Resulta que deseamos obtener solo ese \'atomo. Para ello tenemos que tomar en cuenta que las cadenas numeran sus caracteres del 0 en adelante. C\'omo as\'i? En este caso, por ejemplo, tendr\'iamos que cada letra se puede representar con el sub\'indice que vemos aqu\'i: $\underset{0}{C} \underset{1}{C} \underset{2}{O} \underset{3}{C} \underset{4}{C}$. Entonces, para referirnos al ox\'igeno en la f\'ormula, lo hacemos de la siguiente manera \inlinecode{dietileter[2]}\\

Otra cosa que podemos hacer es referirnos a segmentos de una mol\'ecula. Si deseamos mostrar todo lo que est\'a despu\'es del ox\'igeno de la mol\'ecula anterior, podemos escribir \inlinecode{dietileter[3:]} O si deseamos lo anterior al ox\'igeno \inlinecode{dietileter[:2]} Notemos ahora que en esto \'ultimo, incluimos al ox\'igeno en nuestra selecci\'on, pero al ejecutar el comando, este no se muestra en el resultado. Esto se debe a que Python no considera al \'ultimo elemento al que nos referimos. Como un \'ultimo ejercicio, vamos a mostrar a un azufre en vez de un ox\'igeno utilizando solo lo que hemos visto hasta ahora: \inlinecode{print( dietileter[:2] + "S"\ + dietileter[3:] )		}\\

Antes de proseguir, solo vamos a aclarar algo. Esta forma de representar mol\'eculas a trav\'es de letras es un formato internacional llamado SMILES (\textbf{S}implified \textbf{M}olecular-\textbf{I}nput \textbf{L}ine-\textbf{E}ntry \textbf{S}ystem). Los dobles enlaces se representan con un signo \inlinecode{=}, los triples enlaces con \inlinecode{\#} y para formar ciclos, se numeran los \'atmos que cerrar\'an ese ciclo. Un benceno se ver\'ia entonces de la siguiente manera: \inlinecode{benceno = "\ \hspace{-2mm}C1=CC=CC=C1"}\\

Si deseamos generar nuestros propios SMILES para trabajar m\'as adelante con ellos, podemos visitar el siguiente enlace: \href{https://marvinjs-demo.chemaxon.com/latest/demo.html}{MarvinJS}. Aqu\'i, lo \'unico que necesitamos hacer es dibujar nuestras mol\'eculas, guardarlas y, en ese momento, escoger que el formato de salida sea SMILES. Claro! MarvinJS tambi\'en puede leer formatos smiles, pero con eso trabajaremos m\'as adelante.

\subsubsection{Listas}
Despu\'es de las cadenas, las listas son otro tipo de datos que nos permitir\'a hacer cosas muy interesantes. Las listas nos permiten guardar una serie de datos dentro de ellas. Son como las celdas en las que colocamos informaci\'on en Microsoft Excel, por ejemplo. Otra forma de verlo es como una tira c\'omica. Dentro de cada cuadro se halla una parte importante de la historia: informaci\'on.\\

Las listas en Python se representan como \inlinecode{list}. Para crear una lista colocamos diferentes piezas de informaci\'on separadas por comas dentro de dos corchetes. Supongamos que tenemos los datos de 5 repeticiones de una titulaci\'on \'acido/base. Para crear una lista con esta informaci\'on, vamos a escribir lo siguiente:

\begin{Code}
repeticiones = [22.3, 22.2, 22.4, 22.3, 22.1]
\end{Code}

Por supuesto, no solo se tiene que tratar de n\'umeros. Podemos colocar cadenas a la par de n\'umeros enteros, complejos a la par de flotantes y hasta otras listas dentro de nuestra lista. Como ejercicio, intentemos colocar un dato de cada tipo de los que hemos visto en nuestra lista. Veamos qu\'e sucede y finalmente, veamos lo que hizo nuestro compa\~nero de al lado y lo que result\'o.\\

Volviendo al ejemplo de una titulaci\'on, nos dicen que el cuarto dato estaba mal ingresado. Contrario a lo que pasa con las cadenas, los valores en una lista \emph{s\'i} se pueden cambiar individualmente. C\'omo hacemos esto? Pues es bastante sencillo. Las listas, al igual que las cadenas se accesan as\'i: \inlinecode{repeticiones[3]} Debemos tomar en cuenta que como en nuestro ordenador todas las cosas comienzan con 0, el cuarto dato se halla entonces en la posici\'on 3. El dato que quieren que coloquemos en vez de \inlinecode{22.3} es \inlinecode{22.4}. Entonces, la operaci\'on para cambiar el cuarto dato se ver\'ia de la siguiente manera: \inlinecode{repeticiones[3] = 22.4} Eso fue sencillo en comparaci\'on a cambiar el ox\'igeno por azufre en una cadena. En este caso, solo cambiamos el dato directamente.\\

Otro par de similitudes entre las listas y las cadenas es que se puede obtener su tama\~no de la misma manera: \inlinecode{len(repeticiones)} Tambi\'en podemos mostrar partes de la lista: \inlinecode{repeticiones[2:4]} Pero una cosa que cambia un poco es el agregado y la eliminaci\'on de alg\'un dato. S\'i se pueden agregar datos si sumamos dos listas. Sin embargo, para agregar datos en una lista, hay otras dos formas que nos facilitan esto. Veamos cada caso:

\begin{Code}
\# Sumando listas\\
repeticiones + [25.5]

\# Agregando datos al final\\
repeticiones.append(25.5)

\# Agregando datos en alguna parte en particular\\
repeticiones.insert(2, 25.5)
\end{Code}

Intentemos hacer esto con nuestra lista y veamos qu\'e es lo que pasa en cada caso. Al final, discutamos con nuestro compa\~nero de al lado lo que creemos que sucedio en cada caso.\\

Por supuesto, despu\'es de haberle agregado datos falsos a nuestra lista, queremos quit\'arselos. Despu\'es de todo, no tiene mucho sentido tener datos falsos en un an\'alisis. Para eso, tenemos 2 maneras de hacerlo. La primera quita el dato en la posici\'on que indicamos. El segundo quita la primera vez que aparece el dato que indicamos. Veamos el ejemplo:

\begin{Code}
\# Quitando el dato por su posicion\\
repeticiones.pop(2)

\# Quitando el dato que indicamos\\
repeticiones.remove(25.5)
\end{Code}

Nuevamente, hagamos esto con nuestra lista y veamos qu\'e pasa. Finalmente discutamos con nuestro compa\~nero de al lado a ver qu\'e sucedi\'o con nuestra lista.\\

Otro par de operaciones interesantes que podemos hacer con las listas es ordenarlas o invertir su orden. Veamos esto \'ultimo. Si tomamos nuestra lista y pensamos en tomar solo un segmento, lo hacemos con los corchetes y dos puntos. Pero qu\'e pasar\'a si escribimos esto?

\begin{Code}
\# Invirtiendo el orden de la lista\\
repeticiones[::-1]
\end{Code}

Qu\'e result\'o de hacer esto? Por qu\'e crees que sucedi\'o eso? Qu\'e pasar\'ia si en vez de un \inlinecode{-1} escribi\'eramos otro n\'umero? Intent\'emoslo y veamos qu\'e pasa.\\

Finalmente, vamos a aprender dos comandos. Aunque los dos se parecen y operan de manera similar, el resultado ser\'a un tanto distinto. Vamos a crear dos listas iguales y luego vamos a ordenarlas.

\begin{Code}
primera = [1, 17, 3, 13, 7, 11, 5, 2]\\
segunda = [1, 17, 3, 13, 7, 11, 5, 2]\\
print( primera )\\
primera.sort() \# Ordenamos la lista simplemente\\
print( primera )\\
print( segunda )\\
sorted( segunda ) \# Ordenamos una copia de la lista y la devolvemos\\
print( segunda )
\end{Code}

Si ponemos atenci\'on, fuimos viendo qu\'e suced\'ia con cada lista al ir aplicando dos operaciones distintas. Cu\'al fue el cambio? En qu\'e ocasi\'on podr\'iamos hacer uso de cada una de las dos? Comenta con tu compa\~nero de al lado sobre los resultados.\\

M\'as adelante vamos a ver que las listas nos van a servir para mucho m\'as que solo ver los datos en una fila. Podremos extraer mucha m\'as informaci\'on de ellos cuando aprendamos otro par de \emph{trucos}.

\subsubsection{Diccionarios}
El \'ultimo tipo de dato que vamos a ver son los \emph{diccionarios}. Ahora, estos no son como los diccionarios que conocemos hasta ahora. El concepto, sin embargo, es similar: utilizamos una palabra para hallar informaci\'on sobre ella. Poni\'endonos m\'as t\'ecnicos, un diccionario en Python es igual a una lista, pero en vez de n\'umeros para referirnos a la posici\'on de los datos, utilizamos cadenas para hacerlo. En otras palabras, podemos ponerle nombre a cada casilla. Veamos, pues, c\'omo hacer esto.\\

En Python, los diccionarios se representan como \inlinecode{dict}. Y como este es un taller de Computaci\'on Cient\'ifica, intentaremos construir el diccionario m\'as utilizado por nosotros en nuestra carrera. Los diccionarios se construyen con llaves en vez de corchetes a diferencia de las listas. Adem\'as de eso, para colocarle nombre a cada casilla, incluimos una cadena con el nombre antes del valor, seguido de dos puntos.

\begin{Code}
tabla = $\lbrace$ "H": 1.0079, "He": 4.0026, "Li": 6.94, "Be": 9.0122, "B": 10.81, "\ \hspace{-2mm}C": 12.011, "N": 14.007, "\ \hspace{-2mm}O": 15.9994, "F": 18.998 $\rbrace$
\end{Code}

Nos quedaremos con los primeros 9 elementos, ya que incluir los 112 es un poco tedioso. Nos disponemos ahora a ver c\'omo tomar los datos de nuestro diccionario. Para ello vamos a buscar el peso at\'omico del carbono: \inlinecode{tabla["\ \hspace{-2mm}C"]} Ahora nos damos cuenta de que en vez de buscar una entrada en particular en el diccionario, vamos a buscar una casilla por su nombre. Esto nos facilita almacenar m\'as informaci\'on! No solo estamos guardando el s\'imbolo del elemento, sino su peso at\'omico, y adem\'as somos capaces de buscar sin necesidad de recordar en qu\'e orden estaba todo almacenado.\\

Siguiendo la manera en que trabajamos con las listas, vamos a ver c\'omo agregar datos y quitarlos en los diccionarios. Vale la pena decir que quitar datos se hace de la misma manera que con las listas: con \inlinecode{pop()}. No obstante, en vez de utilizar un n\'umero indicando la casilla, utilizamos el nombre de la casilla. Para agregar datos dentro del diccionario, lo \'unico que debemos hacer es referirnos a una casilla como si ya estuviera creada. En este caso agregaremos el ne\'on: \inlinecode{tabla["Ne"] = 20.180} Eso es todo! Como los diccionarios no tienen un orden, no importa d\'onde coloquemos el dato ni de d\'onde lo quitemos.\\

Finalmente, los diccionarios tienen una caracter\'istica que vale la pena mencionar: si deseamos solo las ``llaves'' del diccionario (es decir, lo nombres de cada casilla), podemos pedirlas as\'i: \inlinecode{tabla.keys()} Esto nos devuelve una lista con los nombres de cada casilla (por si se nos olvidan).\\

Los diccionarios son un tipo de dato muy particular de Python que, si los aprendemos a usar bien, pueden ayudarnos mucho. No los olvidemos e intentemos aprovechar de esta ventaja, porque estos nos pueden facilitar muchas otras cosas y ahorrarnos tiempo en el futuro.

\subsection{Ingresando y Mostrando Datos}
Hemos aprendido qu\'e tipos de datos existen en Python y qu\'e se puede hacer con ellos. Sin embargo, todo esto lo hemos visto probando directamente en la l\'inea de comando. Cuando pensamos en programas, son cosas independientes que corren al click de un \'icono y funcionan de inmediato. Esa es nuestra siguiente meta en nuestro aprendizaje.\\

Lo que nos toca ahorita es ver c\'omo hacer que nuestro ordenador corra varios comandos seguidos, nos pida datos, nos muestre resultados y haga operaciones sin necesidad de que nosotros le ingresemos cada cosa a la vez. Tambi\'en veremos c\'omo hacer que nuestro ordenador lea los datos de un documento, que los interprete y que haga algo con ellos. Finalmente, escribir datos en un documento tambi\'en puede ser una opci\'on cuando se trata de mostrar resultados y guardarlos a la vez. La idea es que nuestro ordenador comience a hacer cosas por nosotros sin que se lo estemos diciendo siempre y sin estarlo supervisando en cada paso.

\subsubsection{Por Parte del Usuario}
Ahora que ya sabemos con qu\'e cosas se pueden trabajar, veamos c\'omo comenzar a hacer una rutina o \textit{script}. Para esto vamos a salirnos de Python en la l\'inea de comando ingresando \inlinecode{quit()}. Dejando la l\'inea de comando abierta, vamos a crear un nuevo directorio \textit{Scripts} dentro de nuestro directorio de documentos. Luego vamos a abrir un editor de texto y vamos a guardar el documento en blanco como \textit{bienvenida.py} en el directorio reci\'en creado.\\

Consideremos ahora que vamos a escribir una rutina. Lo primero que debemos preguntarnos es ``C\'omo se supone que funciona una rutina?'' O en otras palabras, debemos preguntarnos c\'omo la vamos a hacer funcionar despu\'es de haberla escrito. Para hacer que nuestra rutina o script funcione, debemos de decirle a Python que corra el documento en donde se halla lo que escribimos. Esto lo hacemos escribiendo \inlinecode{python bienvenida.py} Claro, ahorita no va a pasar nada si lo hacemos, porque no hay nada en el documento, pero m\'as adelante, esta es la manera de ir viendo lo que hace el c\'odigo que hemos escrito. Pasemos entonces a crear nuestra primera rutina.\\

\noindent En el editor de texto, vamos a escribir lo siguiente:

\begin{Code}
print("Hola mundo!")\\
nombre = raw\_input("\ \hspace{-2mm}Ingresa, por favor, tu nombre: ")\\
print("Hola", nombre + "\ \hspace{-2mm}! Como estas?")
\end{Code}

Intentemos correr nuestra peque\~na rutina y ver qu\'e pasa. Despu\'es de ver el resultado, comenta con tu compa\~nero de al lado: Qu\'e crees que pas\'o? Qu\'e cosas de las que hemos visto antes acabas de utilizar? Qu\'e hace la funci\'on \inlinecode{raw\_input()}? Qu\'e cosa nueva acabamos de aprender a hacer? Tomemos nota de esto y sigamos.\\

Otra funci\'on que nos permite hacer algo similar a \inlinecode{raw\_input()} es solamente \inlinecode{input()}. Estas dos funciones tienen una diferencia en particular. \inlinecode{raw\_input()} toma cualquier cosa que ingresemos y lo convierte en una cadena. No importa si ingresamos n\'umeros, listas, etc. todo ser\'a una cadena. En el caso de \inlinecode{input()}, este solo admite n\'umeros enteros. Cualquier otra cosa que le ingresemos nos dar\'a un error. Sobre esto \'ultimo hablaremos m\'as adelante. Por ahora, vamos a ver c\'omo hacer algo interesante con lo que hemos aprendido.\\

Cerremos nuestro documento \textit{bienvenida.py} y abramos uno nuevo. Esta vez vamos a hacer una rutina que haga algo productivo. Vamos a crear un nuevo documento dentro de nuestro directorio \textit{Scripts}, al que llamaremos \textit{formula.py}. Nuestra rutina devolver\'a la f\'ormula molecular de un compuesto en SMILES. Adem\'as de esto nos devolver\'a tambi\'en el peso molecular. Para ello vamos a ir paso a paso. Lo primero que necesitamos hacer es solicitarle la mol\'ecula en formato SMILES al usuario. Para ello escribimos:

\begin{Code}
molecula = raw\_input("\ \hspace{-2mm}Por favor ingresa la molecula en formato SMILES: ")
\end{Code}

\noindent Luego, vamos a necesitar los pesos at\'omicos de cada elemento para posteriormente poder calcular el peso molecular. Tomando, entonces, el diccionario que hab\'iamos visto la secci\'on anterior, colocamos la segunda l\'inea de c\'odigo en nuestra rutina.

\begin{Code}
molecula = raw\_input("\ \hspace{-2mm}Por favor ingresa la molecula en formato SMILES: ")\\
tabla = $\lbrace$ "H": 1.0079, "He": 4.0026, "Li": 6.94, "Be": 9.0122, "B": 10.81, "\ \hspace{-2mm}C": 12.011, "N": 14.007, "\ \hspace{-2mm}O": 15.9994, "F": 18.998, "Ne": 20.180 $\rbrace$
\end{Code}

\noindent Ahora, considerando que en qu\'imica org\'anica com\'unmente se utiliza los elementos \emph{H}, \emph{C}, \emph{N}, \emph{O} y \emph{F}, vamos a cuantificar solo estos. Esto lo podemos hacer solicit\'andole a nuestro ordenador que cuente cu\'antas veces aparece cierto elemento en el SMILES.

\begin{Code}
molecula = raw\_input("\ \hspace{-2mm}Por favor ingresa la molecula en formato SMILES: ")\\
tabla = $\lbrace$ "H": 1.0079, "He": 4.0026, "Li": 6.94, "Be": 9.0122, "B": 10.81, "\ \hspace{-2mm}C": 12.011, "N": 14.007, "\ \hspace{-2mm}O": 15.9994, "F": 18.998, "Ne": 20.180 $\rbrace$\\
c = molecula.count("\ \hspace{-2mm}C")\\
n = molecula.count("\ \hspace{-2mm}N")\\
o = molecula.count("\ \hspace{-2mm}O")\\
f = molecula.count("\ \hspace{-2mm}F")
\end{Code}

Solo nos queda el hidr\'ogeno. El problema con este es que no aparece de manera expl\'icita en un SMILES, por lo que debemos calcularlo. La \emph{f\'ormula} para calcular los hidr\'ogenos en un compuesto org\'anico es la siguiente:

\begin{equation}
n_H = \left( 2 \cdot n_C + 2 \right) + n_N - n_F - 2 \cdot n_= - 4 \cdot n_\# - 2 \cdot n_{ciclo}
\end{equation}

Si nos damos cuenta, $n_X$ es el n\'umero de \'atomos de $X$. Donde $X$ puede ser carbono, nitr\'ogeno, ox\'igeno, fl\'uor o hidr\'ogeno. El problema son $n_=$, $n_\#$ y $n_{ciclo}$. Ahora, cuantificar los dobles $=$ y triples $\#$ enlaces es sencillo: lo hacemos igual que con los elementos. Pero los ciclos son un poco m\'as complicados. Para eso tenemos que cuantificar cu\'antos n\'umeros hay. Suponiendo que el usuario no ingresar\'a un n\'umero mayor a 5, vamos a ir revisando si los n\'umeros del 1 al 5 se hallan en el SMILES y, de ser as\'i, los agregaremos nuestro conteo.\\

\noindent Nuestro c\'odigo se debe ver de esta manera:

\begin{Code}
molecula = raw\_input("\ \hspace{-2mm}Por favor ingresa la molecula en formato SMILES: ")\\
tabla = $\lbrace$ "H": 1.0079, "He": 4.0026, "Li": 6.94, "Be": 9.0122, "B": 10.81, "\ \hspace{-2mm}C": 12.011, "N": 14.007, "\ \hspace{-2mm}O": 15.9994, "F": 18.998, "Ne": 20.180 $\rbrace$\\
c = molecula.count("\ \hspace{-2mm}C")\\
n = molecula.count("\ \hspace{-2mm}N")\\
o = molecula.count("\ \hspace{-2mm}O")\\
f = molecula.count("\ \hspace{-2mm}F")\\
dobles = molecula.count("\ \hspace{-2mm}=")\\
triples = molecula.count("\ \hspace{-2mm}\#")\\
ciclos = molecula.count("\ \hspace{-2mm}1") / 2 + molecula.count("\ \hspace{-2mm}2") / 2 + molecula.count("\ \hspace{-2mm}3") / 2 + molecula.count("\ \hspace{-2mm}4") / 2 + molecula.count("\ \hspace{-2mm}5") / 2\\
h = (2 * c + 2) + n - f - 2 * dobles - 4 * triples - 2 * ciclos
\end{Code}

\noindent Ya llevamos la mayor parte de la rutina. Ahora solo quedan dos cosas: mostrar la f\'ormula molecular y calcular su peso. Mostrar la f\'ormula es bastante f\'acil ya que tenemos los coeficientes de todos los \'atomos. Para el peso haremos uso de diccionarios. La idea es ir refiri\'endose a cada valor del diccionario y multiplicarlo por el coeficiente que le corresponde. Al final sumaremos esos valores y obtendremos el peso molecular.

\begin{Code}
molecula = raw\_input("\ \hspace{-2mm}Por favor ingresa la molecula en formato SMILES: ")\\
tabla = $\lbrace$ "H": 1.0079, "He": 4.0026, "Li": 6.94, "Be": 9.0122, "B": 10.81, "\ \hspace{-2mm}C": 12.011, "N": 14.007, "\ \hspace{-2mm}O": 15.9994, "F": 18.998, "Ne": 20.180 $\rbrace$\\
c = molecula.count("\ \hspace{-2mm}C")\\
n = molecula.count("\ \hspace{-2mm}N")\\
o = molecula.count("\ \hspace{-2mm}O")\\
f = molecula.count("\ \hspace{-2mm}F")\\
dobles = molecula.count("\ \hspace{-2mm}=")\\
triples = molecula.count("\ \hspace{-2mm}\#")\\
ciclos = molecula.count("\ \hspace{-2mm}1") / 2 + molecula.count("\ \hspace{-2mm}2") / 2 + molecula.count("\ \hspace{-2mm}3") / 2 + molecula.count("\ \hspace{-2mm}4") / 2 + molecula.count("\ \hspace{-2mm}5") / 2\\
h = (2 * c + 2) + n - f - 2 * dobles - 4 * triples - 2 * ciclos\\
C = ( "\ \hspace{-2mm}C"\ + str( c ) ) * bool( c )\\
H = ( "\ \hspace{-2mm}H"\ + str( h ) ) * bool( h )\\
N = ( "\ \hspace{-2mm}N"\ + str( n ) ) * bool( n )\\
O = ( "\ \hspace{-2mm}O"\ + str( o ) ) * bool( o )\\
F = ( "\ \hspace{-2mm}F"\ + str( f ) ) * bool( f )\\
print( "\ \hspace{-2mm}La formula molecular es: "\ + C + H + N + O + F )\\
mC = tabla["\ \hspace{-2mm}C"] * c\\
mH = tabla["\ \hspace{-2mm}H"] * h\\
mN = tabla["\ \hspace{-2mm}N"] * n\\
mO = tabla["\ \hspace{-2mm}O"] * o\\
mF = tabla["\ \hspace{-2mm}F"] * f\\
peso\_molecular = mC + mH + mN + mO + mF\\
print( "\ \hspace{-2mm}El peso molecular es: "\ + str( peso\_molecular ) + "\ \hspace{-2mm}\ g/mol"\ )
\end{Code}

Este es un buen momento para detenerse y reflexionar: Acabamos de \emph{multiplicar} booleanos? De hecho s\'i! Esta es una de esas ventajas de Python. Los valores \inlinecode{True} y \inlinecode{False} son interpretados como $1$ y $0$ respectivamente. Esto hace que mostrar o no mostrar algunas cosas sea f\'acil.\\

Se supone que a estas alturas ya terminamos nuestra rutina. Lo \'unico que queda para confirmar esto es probarla. Para esto es conveniente comenzar con alguna mol\'ecula que no sea tan dif\'icil que rompa lo que acabamos de escribir, pero tampoco algo demasiado f\'acil con lo que no sintamos que nos est\'a ayudando. Guardemos nuestro trabajo y probemos la rutina escribiendo \inlinecode{python formula.py}. Cuando nos sea solicitado, ingresaremos \inlinecode{CNCC(C1=CC(=C(C=C1)O)O)O}; el SMILES de la epinefrina. Funcion\'o? Te arroj\'o los resultados que esperabas? Discute con tu compa\~nero de al lado y toma en cuenta que los resultados de esta mol\'ecula se hallan en su p\'agina de \href{https://pubchem.ncbi.nlm.nih.gov/compound/epinephrine}{PubChem}.\\

Puedes continuar probando con otros SMILES, si as\'i lo deseas. Explora los l\'imites de nuestra rutina y, de hallarlos, explica por qu\'e es que falla en ciertos casos. Por otra parte, piensa en todo lo que podr\'ias hacer de utilizar este tipo de rutinas. Qu\'e otras cosas podr\'ias hacer?

\subsubsection{Desde Documentos}
Continuando con la idea de tomar datos y enviar datos fuera del script, ahora vamos a intentar escribir dos rutinas m\'as peque\~nas que la anterior en las que vamos a crear un documento de texto desde nuestro script, y luego vamos a leer el contenido del mismo documento con otro script. Esto nos puede ser muy \'util a la hora de generar resultados de un c\'alculo largo o para interpretar los datos generados por otro programa. Procedamos con un par de ideas sencillas.\\

Vamos a guardar una mol\'ecula en formato SMILES en un documento de texto (como los que hacemos con el block de notas en windows). Para ello vamos a crear un peque\~no script \textit{escribiendo.py} en nuestro directorio \textit{Scripts}. En \'el vamos a escribir lo siguiente:

\begin{Code}
molecula = raw\_input("\ \hspace{-2mm}Ingresa una molecula en formato SMILES: ")\\
documento = open("molecula.smiles", "w")\\
documento.write(molecula)\\
documento.close()
\end{Code}

Eso es todo! Veamos qu\'e hace cada parte de nuestro nuevo script. Al principio solo le pedimos al usuario que nos de una mol\'ecula (i.e. informaci\'on que podamos guardar en un documento). Luego abrimos un documento de nombre \textit{molecula.smiles} en modo \emph{escritura}: por eso escribimos la \inlinecode{"w"}. Luego, pedimos que a \inlinecode{documento} (la variable que representa nuestro documento nuevo) se le escriba la cadena que tenemos guardada en \inlinecode{molecula} Finalmente, rompemos la conecci\'on entre la variable \inlinecode{documento} y \textit{molecula.smiles} cerrando el documento. De esta forma, la mol\'ecula queda guardada.\\

Antes de continuar, vamos a correr nuestro script y a almacenar alguna mol\'ecula en SMILES. Recordemos lo que hemos aprendido antes sobre este sistema para representar mol\'eculas. Posteriormente veamos que la mol\'ecula se haya guardado e intentemos visualizarla con alg\'un programa como Avogadro o PyMOL. Si no sabemos c\'omo hacer esto, preguntemos a nuestro compa\~nero de al lado o a quien nos est\'a dando el taller.\\

Para continuar, vamos a leer nuestra nueva mol\'ecula y vamos a mostrarla en pantalla. Este script ser\'a todav\'ia m\'as f\'acil que el anterior. Primero vamos a abrir el documento SMILES, luego vamos a leer lo que hay en \'el, vamos a cerrar el documento y luego vamos a mostrar en pantalla lo que hab\'ia all\'i. Veamos el c\'odigo.

\begin{Code}
documento = open("molecula.smiles", "\ \hspace{-2mm}r")\\
informacion = documento.readlines()\\
documento.close()\\
print(informacion)
\end{Code}

Vamos a guardar nuestro script y correrlo. Qu\'e pas\'o? Funcion\'o como esper\'abamos? Comentemos con nuestro compa\~nero de al lado.\\

Despu\'es de eso pensemos un momento: Podr\'iamos hacer que el usuario escoja el nombre del documento a abrirse? Podr\'iamos hacer que el usuario le ponga el nombre al documento? Qu\'e otras cosas podr\'iamos hacer ahora que ya sabemos c\'omo guardar informaci\'on en documentos y leerla? Pensemos en todas las posibilidades ahora y tomemos nota para despu\'es.

\subsubsection{Desde la L\'inea de Comando}
Otra manera de ingresar datos a nuestro ordenador es desde la l\'inea de comando. Ya en casos anteriores hemos hecho esto, pero no nos hemos dado cuenta (e.g. \inlinecode{ls -l}, \inlinecode{wc -w}). Esta vez vamos a leer estos argumentos y vamos a pedirle a Python que nos los entregue para poder utilizarlos. Recordando la \textbf{Sesi\'on 1}, vamos a escribir una rutina que nos calcule la media de un grupo de datos.\\

Vamos a comenzar entonces creando un documento \textit{stats.py} en nuestro directorio \textit{Scripts}. Abrimos el documento y vamos a continuar con un par de comandos nuevos. El primero es \inlinecode{import sys}. Este solo nos est\'a ampliando las funciones de Python. De esta manera vamos a tomar todos los datos que nos dar\'a el usuario. Llevamos entonces lo siguiente:

\begin{Code}
import sys\\
datos = sys.argv
\end{Code}

Ahora ya tenemos los datos que ingresar\'a el usuario en forma de una \emph{lista de cadenas}. Es por ello que tenemos 2 problemas. El primero es que el primer dato en la lista es el nombre de nuestro documento (i.e. \emph{stats.py}). El segundo problema es que los valores ingresados son \emph{cadenas}. El primer problema es bastante f\'acil de solucionar. Se te ocurre c\'omo?\footnote{Recuerda que las listas se pueden partir en pedazos.} El segundo problema requiere el segundo comando nuevo que vamos a ver. Este nos permitir\'a convertir todas las cadenas de la lista en n\'umeros. Veamos c\'omo hacer esto:

\begin{Code}
import sys\\
datos = sys.argv\\
datos = datos[1:]\\
datos = map( float, datos)
\end{Code}

Observa bien la \'ultima l\'inea y comenta con tu compa\~nero de al lado. Qu\'e crees que hace \inlinecode{map()}? Continuando, vemos que ya tenemos todos nuestros datos en una lista. Lo que nos queda ahora es sumarlos todos, dividirlos entre el n\'umero de datos que tenemos y, claro, mostrarlos al usuario. Vamos por pasos, primero la suma, luego contar cuantos elementos (\textit{items}) hay en la lista, posteriormente dividir y finalmente mostrar los datos. Escribamos esa parte:

\begin{Code}
import sys\\
datos = sys.argv\\
datos = datos[1:]\\
datos = map( float, datos)\\
suma = sum( datos )\\
items = len( datos )\\
print( suma / items )
\end{Code}

Muy bien, ya terminamos. Ahora, para correr la rutina solo debemos llamar a \inlinecode{python stats.py}, pero al final, debemos agregar la lista de n\'umeros de los que queremos calcular la media. Intentemos con los resultados de un an\'alisis:  \inlinecode{python stats.py 15.45 14.93 15.10 15.32 14.99 15.04}. Qu\'e resultado nos dio? Revisemos con nuestros compa\~neros y veamos si fue igual.\\

Pensemos ahora en lo que podr\'iamos hacer con este tipo de rutina. Qu\'e pasar\'ia si combin\'aramos esto con lo que aprendimos en las secciones pasadas? Que una rutina abra un documento que nosotros especifiquemos en la l\'inea de comando y calcule las propiedades de un SMILES? Podr\'ia ser. Pero no nos quedemos con eso. Qu\'e otras utilidades le podr\'iamos dar es esto en \emph{nuestros} propios ambientes y carreras? Tomemos nota y continuemos.

\pagebreak

\subsection{Errores}
Esta es una de las partes que a muchas personas no les gusta. Cuando de repente algo no funciona como se esperaba y resultamos con una de dos: mucho texto incomprensible de lo que solo interpretamos que cometimos un error, o que el resultado que esper\'abamos no aparezca. Ambos casos suelen desanimarnos o molestarnos mucho. Pero en este caso, como estamos aprendiendo y la idea es cometer errores, podemos sentirnos tranquilos si algo as\'i sucede. Este tema es, de hecho, hasta evitado a veces en cursos o talleres porque se llega a creer que no vale la pena ver errores. Sin embargo, un aspecto clave de hacer cosas a nivel anal\'itico y de calidad en Computaci\'on Cient\'ifica es saber qu\'e error se cometi\'o, en d\'onde y por qu\'e. Finalmente, debemos ser capaces de corregirlo, optimizar nuestro script y hacer que este sea lo mejor que se pueda.\\

En esta secci\'on vamos a ver la anatom\'ia de un error en Python, qu\'e errores podemos esperar cometer y c\'omo hallar aquellos que no rompen nuestro script, pero tampoco dicen que haya un error en alguna parte. Pero vamos por pasos. Lo primero que debemos entender es que nuestro ordenador, a pesar de ser muy capaz, no posee inteligencia como tal. Este solo hace lo que nosotros le pedimos hacer y no o es capaz de proponer algo o pensar en soluciones a problemas. Esto nos lleva a que necesitamos varias reglas generales para poder indicarle qu\'e hacer; este no sabr\'a qu\'e es lo que deseamos. Estas reglas son una sintaxis que no podemos dejar de utilizar. Si lo hacemos, el ordenador no har\'a lo que nosotros queremos que haga y nos frustraremos.\\

Para este punto del taller, es posible y muy probable que ya hayamos cometido alg\'un error o que hayamos visto uno al menos. Pero para no recurrir a posibles casos pasados, veamos uno de verdad. En una l\'inea de comando con Python abierto ingresemos: \inlinecode{print("Hola mundo"\ \hspace{-2mm}!)} Resulta que corremos eso porque se nos pas\'o el orden de las cosas por teclearlas r\'apido, y el resultado es este: \inlinecode{\color{red} SyntaxError: invalid syntax} Y pues, claro! Colocamos un caracter en donde no iba. Pero nada se ha arruinado, todo sigue funcionando en nuestro ordenador y el sol sigue brillando afuera. Qu\'e sucedi\'o entonces? Realmente es un sencillo. Python se detuvo donde hall\'o un error de sintaxis, advirti\'o qu\'e tipo de error era al mostr\'arnoslo y se qued\'o all\'i. Nos toca arreglarlo. Pero qu\'e pasa cuando el error ya es un script y no en la l\'inea de comando? Veamos esos casos.\\

Vamos a ejecutar un script que tenemos en nuestro directorio de documentos. Este se llama \textit{errores.py} y nos va a ir mostrando paso a paso los tipos de error con los que nos podemos ir topando\footnote{Solo vamos a ver los errores m\'as comunes que podemos cometer. Hay muchos otros tipos de error, pero no los cubriremos aqu\'i.}. Adem\'as de ejecutar el script, lo vamos a abrir con un editor de texto y ver qu\'e pasa en cada parte. Vamos a ir uno por uno entendiendo los errores.\\

\subsubsection{Error de Sintaxis}
Este ser\'a el primer error que aparecer\'a al ejecutar el script. Si nos damos cuenta, es el mismo error que antes. Pero esta vez, a la hora de aparecer en la l\'inea de comando, obtuvimos m\'as informaci\'on que antes. Esta vez apareci\'o lo siguiente:

\begin{Code}
{\color{red} \hspace*{4mm} File "\ \hspace{-2mm}errores.py", line 2\\
\hspace*{8mm} print "\ \hspace{-2mm}Error de\ ".\ "\ \hspace{-2mm}sintaxis."\\
\hspace*{63mm}\^\\
SyntaxError: invalid syntax}
\end{Code}

Vamos a ver l\'inea por l\'ina qu\'e es lo que el error nos est\'a diciendo. En la primera l\'inea nos explica que el error se halla en el script \textit{errores.py} y, dentro de este, en la l\'inea 2. En la segunda l\'inea se nos muestra la l\'inea de c\'odigo donde sucedi\'o el error. En la tercera se nos intenta dar una idea de d\'onde se halla el error. Esta parte no siempre es precisa; a veces se nos indica realmente donde sucedi\'o el error, y a veces no tanto. En la l\'inea final solo se nos muestra qu\'e tipo de error es.\\

Con respecto de esto \'ultimo, podemos decir que un error de sintaxis es, generalmente, cuando escribimos algo de manera incorrecta. Nuestro ordenador entonces nos advierte sobre esto y nos toca corregirlo. Para continuar, vamos a eliminar las dos l\'ineas del script \textit{errores.py} en donde se hallaba el error de sintaxis, vamos a guardarlo y vamos a volverlo a correr.

\subsubsection{Error de Nombre}
Esta vez obtuvimos un error diferente. Este se nos dice que es un error de nombre cuando leemos la \'ultima l\'inea. De hecho, se nos dice de una vez cu\'al fue el error: hay una variable que no est\'a definida. Al revisar el resto de la informaci\'on sobre el error, nos damos cuenta de que este cambi\'o un poco con respecto del anterior.

\begin{Code}
{\color{red}Traceback (most recent call last):\\
\hspace*{4mm} File "\ \hspace{-2mm}errores.py", line 2, in <module>\\
\hspace*{8mm} print "mi\_variable"\\
NameError: name `mi\_variable' is not defined}
\end{Code}

La primera l\'inea solo nos avisa que Python no solo se ha topado con el error, sino que lo ha buscado. Esto solo significa que cuando hall\'o que el script se refiri\'o a una variable, esta fue \emph{buscada} por todo el script y como no se hall\'o, hubo un error. Lo dem\'as ya lo hemos visto antes: indicaciones sobre el script y la l\'inea donde se halla el error y la muestra de esa l\'inea.\\

Entonces, en resumen, podemos decir que un error de nombre se da generalmente cuando nos referimos a algo que no existe. Para continuar borremos esas l\'ineas y volvamos a correr el script.

\subsubsection{Error de Divisi\'on entre Cero}
El siguiente error que vemos es algo que desde nuestros primeros cursos de matem\'atica en el colegio nos advirtieron que no se pod\'ia hacer: Uno no puede dividir entre cero; eso est\'a prohibido. Este error se parece mucho m\'as al anterior, y es que Python intent\'o buscarle sentido a lo que hicimos, pero no se lo hall\'o. Veamos el error.

\begin{Code}
{\color{red}Traceback (most recent call last):\\
\hspace*{4mm} File "\ \hspace{-2mm}errores.py", line 3, in <module>\\
\hspace*{8mm} otro\_numero = algun\_numero / 0\\
ZeroDivisionError: integer division or modulo by zero}
\end{Code}

Vemos que hubo una b\'usqueda, se nos indica el script, la l\'inea y se nos muestra la l\'inea de c\'odigo con el error. Finalmente, se nos advierte qu\'e tipo de error es y lo que significa. Este tipo de error nos aparecer\'a al dividir entre 0 o al intentar calcular el m\'odulo de un n\'umero entre 0. Eliminemos las l\'ineas necesarias del script, volv\'amoslo a correr y sigamos avanzando.

\subsubsection{Error de Tipo}

\begin{Code}
{\color{red}Traceback (most recent call last):\\
\hspace*{4mm} File "\ \hspace{-2mm}errores.py", line 4, in <module>\\
\hspace*{8mm} resultado = lista\_a * lista\_b\\
TypeError: can't multiply sequence by non-int of type `list'}
\end{Code}

En este error nos damos cuenta de que se trata de tipos. Y claro, un error de tipo se da cuando intentamos alguna operaci\'on sobre tipos de datos que no la permiten. La anatom\'ia del error es la misma que antes, solo que la \'ultim l\'inea nos explica qu\'e es lo que pasa. De nuevo, eliminemos las l\'ineas necesarias para quitar este error del script, volv\'amoslo a correr y continuemos.

\subsubsection{Error de Atributo}
En el caso de estos errores ya solo vamos a enfocarnos en la \'ultima l\'inea. Como nos hemos dado cuenta, ya entendimos qu\'e es lo que tienen las dem\'as y la informaci\'on que nos dan.

\begin{Code}
{\color{red} AttributeError: `str' object has no attribute `append'}
\end{Code}

Este error nos aparece cuando intentamos utilizar alguna funci\'on de un objeto o tipo de dato que no posee tal funci\'on. En este caso, intentamos agregarle un elemento a una cadena como lo har\'iamos a una lista. Esto, como vemos, no se puede hacer. Eliminemos lo necesario del script y volvamos a correr.

\subsubsection{Error de Valor}

\begin{Code}
{\color{red} ValueError: could not convert string to float: esto es texto}
\end{Code}

Si observamos bien el script, est\'abamos intentando transformar una cadena con texto adentro en un n\'umero. Eso no se puede hacer. Y por eso se nos advierte que hay un error de valor. Este tipo de error generalmente aparece cuando intentamos alguna transformaci\'on de un tipo de datos en otro que no es permitido. Continuemos eliminando l\'ineas y volviendo a correr.

\subsubsection{Error de \'Indice}

\begin{Code}
{\color{red} IndexError: list index out of range}
\end{Code}

Este es uno de los errores m\'as comunes que vamos a cometer. Sucede cuando intentamos referirnos a un \'indice inexistente de una lista. En otras palabras, enviamos a nuestro ordenador a buscar una direcci\'on que no existe. Es muy f\'acil de cometer, porque generalmente (como veremos m\'as adelante) vamos cambiando la cantidad de datos en una lista. Una idea para no cometer este error es siempre revisar el tama\~no de la lista antes de accesarla. Continuemos igual que antes para llegar al \'ultimo tipo de error que vamos a ver.

\subsubsection{Error de Entrada y Salida de Datos (IO)}

\begin{Code}
{\color{red} IOError: [Errno 2] No such file or directory: `algun\_documento.txt'}
\end{Code}

Ya hab\'iamos visto de qu\'e se trata la entrada y salida de datos: leer informaci\'on de alguna parte y enviarla/guardarla. En este caso, intentamos abrir un documento que no existe buscando leer de \'el. Nuestro ordenador inmediatamente nos advierte que ese documento no existe con un error de IO (\textbf{I}nput - \textbf{O}utput).

\subsubsection{Quitando Insectos}
Cuando corremos nuestro script y resulta que no hay errores, pero no obtenemos resultado o el resultado no es lo que esper\'abamos, existe una manera hallar el problema. Este proceso se llama \emph{debuggeo} (o \emph{debugging} en ingl\'es), que significa quitar insectos.\\

Hay ambientes de desarrollo integrados (IDEs) que ya hacen esto por nosotros. Sin embargo, en el caso de no contar con uno, o no querer usarlo porque para lo que estamos haciendo un IDE es mucho, la manera m\'as sencilla de hallar d\'onde se halla el problema, es imprimiendo el valor de las variables importantes (o todas, de ser necesario) en los pasos cr\'iticos del script. As\'i podremos ver qu\'e valor est\'a tomando cada una y hallar si todas se est\'an comportando como se debe. De hallar una que no tenga el valor deseado, sabremos que el error est\'a antes de la l\'inea en la que se imprimi\'o el valor de la variable. As\'i podemos revisar el c\'odigo que se halla antes y corregir el error.

\subsubsection{Capturando Errores}
Algunas veces no vamos a saber si una parte de nuestro c\'odigo va a generar un error. Esto suele pasar cuando dependemos de que un usuario interact\'ue con nuestro programa o script. Para evitar que nuestro script se rompa y nos veamos en un aprieto, existe una forma de que nuestro ordenador \emph{pruebe} ejecutar un comando, y si este da un error, tome otra v\'ia como alternativa.\\

Supongamos que necesitamos que un usuario ingrese un n\'umero entero. Para ello utilizamos el siguiente c\'odigo: \inlinecode{opcion = input("Por favor ingrese un numero del 1 al 10: ")} El usuario puede no hacer caso e ingresar un n\'umero con decimales, o peor a\'un, texto. Esto nos generar\'ia un error. Por esta raz\'on decidimos hacer una prueba.

\begin{Code}
try:\\
\hspace*{4mm} opcion = input("Por favor ingrese un numero del 1 al 10: ")\\
except \textit{Excpetion} as e:\\
\hspace*{4mm} print "Se ha producido un error:", e\\
\hspace*{4mm} print "Ha ingresado una opci\'on no-valida."
\end{Code}

Si nos damos cuenta, capturamos el error (si es que se produce uno) y almacenamos sus caracter\'isticas en la variable \inlinecode{e}. De no producirse el error, el c\'odigo solo se seguir\'a ejecutando sin entrar al \inlinecode{except}. Algo que vale la pena mencionar es sobre la palabra \inlinecode{\textit{Exception}}. Esta palabra nos permite capturar cualquier error. Pero si deseamos capturar un error en particular, podemos colocar alguno de los tipos de error que hab\'iamos estado viendo hasta ahora. Eso s\'i, si el error producido en el \inlinecode{try:} es diferente al que se est\'a esperando en el \inlinecode{except}, el error no ser\'a capturado y nuestro script se romper\'a. E.g.

\begin{Code}
try:\\
\hspace*{4mm} 1 / 0\\
except \textit{IOError} as e:\\
\hspace*{4mm} print "Se ha producido un error:", e
\end{Code}

Esto, por seguro, producir\'a un error que no ser\'a capturado. El error esperado es de entrada o salida de datos, y el error producido es de divisi\'on entre cero. Por esta raz\'on, hay que tener mucho cuidado con el tipo de error que se est\'a esperando\footnote{Y para aquellos que piensen en usar \textit{Exception} siempre, consideren que depende del tipo de error es la medida que podemos tomar despu\'es en nuestro script.}.

\subsection{Ejercicio Final: Temperaturas}
Como complemento de lo que ya hemos visto hasta ahora, nos queda un peque\~no ejercicio. Al realizar este vamos a comprender un poco m\'as sobre la l\'ogica utilizada al construir programas. Como nuestras herramientas son limitadas, este ejercicio se enfoca m\'as en hacer uso inteligente de nuestros recursos. Comencemos, entonces, buscando condensar la mayor parte de esta sesi\'on.

Para comenzar, vamos a crear un documento llamado \textit{temperaturas.py} en nuestro directorio de \textit{Scripts}. En \'el vamos a intentar hacer una peque\~na rutina en la que ingresaremos una temperatura y esta nos ser\'a devuelta en otras unidades. Veamos, pues, c\'omo hacer esto.\\

\noindent Lo primero que haremos ser\'a dar la bienvenida a nuestro script. Para eso ingresaremos algo como esto (el texto lo podemos cambiar si as\'i lo deseamos):

\begin{Code}
print("********** Convertidor **********")
\end{Code}

\noindent Luego, deseamos dar las opciones de qu\'e tipo de conversi\'on se puede hacer, as\'i que seguimos mostrando cosas en pantalla:

\begin{Code}
print("********** Convertidor **********")\\
print("1. F -> C \textbackslash t  2. C -> F")\\
print("3. C -> K \textbackslash t  5. K -> F")\\
print("7. F -> K \textbackslash t 11. K -> F")
\end{Code}

\noindent Ahora pediremos al usuario ingresar una opci\'on para calcular la temperatura y luego, la cifra que desea calcular:

\begin{Code}
print("********** Convertidor **********")\\
print("1. F -> C \textbackslash t  2. C -> F")\\
print("3. C -> K \textbackslash t  5. K -> F")\\
print("7. F -> K \textbackslash t 11. K -> F")\\
opcion = input("\ \hspace{-2mm}Ingrese el numero de operacion que desea realizar: ")\\
temperatura = raw\_input("\ \hspace{-2mm}Ingrese la temperatura a ser calculada: ")
\end{Code}

\noindent Ya con esto nos queda transformar la temperatura en un n\'umero y hacer el c\'alculo. Para ello vamos a hacer un truco interesante. Como no conocemos todav\'ia ninguna manera en la que podamos escoger alguna de las operaciones vamos a calcular todas, pero vamos a devolver solo la que nos interesa. C\'omo? Con dos operaciones matem\'aticas. Veamos c\'omo se ve nuestro c\'odigo ya terminado:

\begin{Code}
print("********** Convertidor **********")\\
print("2. F -> C \textbackslash t  3. C -> F")\\
print("5. C -> K \textbackslash t  7. K -> F")\\
print("11. F -> K \textbackslash t 13. K -> F")\\
opcion = input("\ \hspace{-2mm}Ingrese el numero de operacion que desea realizar: ")\\
temperatura = raw\_input("\ \hspace{-2mm}Ingrese la temperatura a ser calculada: ")\\
temperatura = float(temperatura)\\
print( (1 - opcion \% 2) * ("Temperatura en grados C: "\ +\ str((temperatura - 32) * 5.0/9)) )\\
print( (1 - opcion \% 3) * ("Temperatura en grados F: "\ +\ str(temperatura * 9.0/5\ +\ 32)) )\\
print( (1 - opcion \% 5) * ("Temperatura en grados K: "\ +\ str(temperatura\ +\ 273.15)) )\\
print( (1 - opcion \% 7) * ("Temperatura en grados C: "\ +\ str(temperatura - 273.15)) )\\
print( (1 - opcion \% 11) * ("Temperatura en grados K: "\ +\ str((temperatura - 32) * 5/9\ +\ 273.15)) )\\
print( (1 - opcion \% 13) * ("Temperatura en grados F: "\ +\ str((temperatura - 273.15) * 9.0/5\ +\ 32)) )
\end{Code}

Veamos qu\'e pasa aqu\'i. Convertimos la temperatura de una cadena a un n\'umero utilizando \inlinecode{float()} Es importante tomar nota de esto! Podemos convertir un tipo de datos en el otro mediante este tipo de operaci\'on. Luego comenzamos a mostrar cosas en pantalla con \inlinecode{print} A continuaci\'on tenemos una operaci\'on algo extra\~na de la que hablaremos al final. El resultado de esta se multiplica por una cadena a la cual estamos sumando otra cadena. S\'i, la segunda parte es una cadena tambi\'en, porque a pesar de que el resultado de la operaci\'on matem\'atica sea un n\'umero, estamos usando la funci\'on \inlinecode{str()} para convertirlo en una cadena. En otras palabras, estamos combinando el resultado con la presentaci\'on de este en una sola cadena. Hasta all\'i vamos bien. Ahora s\'i, por qu\'e estamos multiplicando una cadena con otro valor de una operaci\'on extra\~na?\\

Veamos detenidamente la operaci\'on un momento. Si consideramos que \inlinecode{opcion} puede tomar solo los valores 2, 3, 5, 7, 11 o 13\footnote{Por qu\'e escogimos esos n\'umeros? Leamos hasta el siguiente p\'arrafo y pensemos por qu\'e esos n\'umeros y no otros.}, entonces notamos que las operaciones de m\'odulo van a dar como resultado 0 si los n\'umeros a ambos lados del operador son iguales, o alg\'un n\'umero entero de 1 en adelante cuando se trate de n\'umeros diferentes. Eso significa que \inlinecode{(1 - opcion \% n)} va a ser igual a 1 cuando \inlinecode{opcion} sea igual a \inlinecode{n}. Si no, ser\'a 0 o alg\'un n\'umero negativo. Perfecto! Y esto c\'omo afecta a nuestra cadena de al lado? Pues es sencillo. Si una cadena la multiplicamos por alg\'un entero positivo, se va a repetir ese n\'umero de veces. En el caso de multiplicarla por 0 o alg\'un n\'umero negativo, desaparecer\'a. Esto, ya vi\'endolo en nuestro peque\~no script, se traduce en: Solo se mostrar\'a el resultado de la operaci\'on escogida.\\

Antes de continuar a la siguiente parte, analicemos que en este script solo hemos tomado datos, mostrado datos, calculado con un poco de aritm\'etica y propiedades de cadenas. Pensemos c\'omo podr\'iamos optimizar nuestro script para que no hayan tantas cosas que se repitan una y otra vez dentro de \'el\footnote{Considerando lo que vimos en la secci\'on 2.2, pensemos qu\'e pasar\'ia al convertir muchas cosas en variables.}. Adem\'as, comentemos con nuestro compa\~nero de al lado lo que hemos aprendido y c\'omo entendimos que funcionaba el script. Tomemos nota de nuestros hallazgos y d\'emonos un momento antes de continuar.

\newpage

\subsection{Comentarios Finales}
Muchas felicidades! Has completado el cuarto d\'ia del Taller. Ahora est\'as en la capacidad de crear rutinas peque\~nas, trabajar con lo que te ingrese un usuario, con el contenido de un documento y con datos ingresados desde una l\'inea de comando. Adem\'as sabes identificar errores y usarlos a tu favor, si llegara a ser necesario.\\

La importancia de los peque\~nos scripts o rutinas no debe de ser hecho de menos. Estos han sido la semilla de la que ha nacido la computaci\'on cient\'ifica. De nuevo felicitaciones por haber terminado esta sesi\'on; ya alcanzaste un nuevo nivel en este tema. En la pr\'oxima sesi\'on veremos c\'omo ahorrarnos trabajo con tareas repetitivas, c\'omo hacer uso de otros \emph{peque\~nos scripts} de otras personas que nos ayudar\'an a facilitarnos mucho trabajo. Queda mucho por descubrir.

\section*{Licencia}

\noindent \includegraphics{img/cc_big.png}

\noindent Taller de Computaci\'on Cient\'ifica para Ciencias Qu\'imicas by \href{http://github.com/zronyj/TQCA}{Rony J. Letona} is licensed under a \href{http://creativecommons.org/licenses/by-sa/4.0/}{Creative Commons Attribution-ShareAlike 4.0 International License}.

\end{document}
